%%%%%%%%%%%%%%%%%%%%%%%%%%%%%%%%%%%%%%%%%%%%%%%%%%%%%%%%%%%%%%%%%%%%%%%%%%%%%%%%
% F.CSM308 Машин сургалт — Бие даалт 1
% ШУТИС, МХТС-ийн өгүүллийн загварын дагуу (XeLaTeX)
% Хөрвүүлэх: latexmk -xelatex BD1_DKT_B231910012.tex
%%%%%%%%%%%%%%%%%%%%%%%%%%%%%%%%%%%%%%%%%%%%%%%%%%%%%%%%%%%%%%%%%%%%%%%%%%%%%%%%
\documentclass[10pt,twocolumn]{article}
\usepackage[a4paper,left=1.7cm,right=1.7cm,top=2.0cm,bottom=2.0cm,columnsep=0.6cm]{geometry}
\usepackage{fontspec}
\setmainfont{Times New Roman}
\setmonofont{Courier New}
\newfontfamily\cyrillicfonttt{Courier New}
\usepackage{polyglossia}
\setdefaultlanguage{russian}   % кирилл үгийн таслал (монгол хэлний загвар байхгүй)
\setotherlanguage{english}
\usepackage{amsmath,amssymb,bm}
\usepackage{array,colortbl,booktabs}
\usepackage{xcolor}
\usepackage{titlesec}
\usepackage{caption}
\usepackage{tikz,pgfplots}
\pgfplotsset{compat=1.17}
\usetikzlibrary{positioning,arrows.meta}
\usepackage[hidelinks]{hyperref}


% ---------- Гарчиг, дэд гарчиг ----------
\titleformat{\section}{\centering\normalfont\bfseries}{\Roman{section}.}{0.45em}{\MakeUppercase}
\titlespacing*{\section}{0pt}{10pt}{5pt}
\titleformat{\subsection}{\normalfont\itshape}{\Alph{subsection}.}{0.45em}{}
\titlespacing*{\subsection}{0pt}{6pt}{3pt}
\renewcommand{\thesection}{\Roman{section}}
\renewcommand{\thesubsection}{\thesection-\Alph{subsection}}

% ---------- Зураг, хүснэгтийн тайлбар ----------
\DeclareCaptionLabelFormat{mnfig}{#2-р зураг}
\captionsetup[figure]{labelformat=mnfig,labelsep=period,font={small,bf,it},justification=centering}
\newcommand{\tabtitle}[1]{\refstepcounter{table}%
  {\centering\small\MakeUppercase{#1}\par}%
  {\raggedleft\small\itshape\thetable-Р ХҮСНЭГТ.\par}\vspace{2pt}}
\renewcommand{\arraystretch}{1.12}
\newcolumntype{L}[1]{>{\raggedright\arraybackslash}p{#1}}

\setlength{\parindent}{1.2em}
\pagestyle{empty}
\sloppy

\newcommand{\x}{\mathbf{x}}
\newcommand{\h}{\mathbf{h}}
\newcommand{\y}{\mathbf{y}}

\begin{document}

\twocolumn[{%
\begin{center}
{\Large\bfseries
СУРАГЧИЙН МЭДЛЭГИЙГ МӨШГИХ ГҮН СУРГАЛТ:\\[1pt]
DEEP KNOWLEDGE TRACING-ИЙН ОНОЛЫН ШИНЖИЛГЭЭ\\[1pt]
БА ЖИШИГ ӨГӨГДЛИЙН АУДИТ\par}
\vspace{12pt}
{\bfseries Бавуудоржийн ӨСӨХБАЯР\par}
\vspace{2pt}
Монгол улс, Улаанбаатар, ШУТИС, Мэдээлэл, холбооны технологийн сургууль, Компьютерын ухааны салбар\par
\end{center}
\vspace{6pt}
\begin{minipage}{\textwidth}
\leftskip=0.5cm \rightskip=0.5cm
\noindent\textbf{\textit{Хураангуй:}} \textbf{Knowledge tracing буюу мэдлэг мөшгих нь сурагчийн
өмнөх оролдлогуудаас дараагийн бодлогыг зөв бодох магадлалыг таамаглах бодлого бөгөөд адаптив
сургалтын системийн суурь юм. Энэ ажилд Piech нарын санал болгосон Deep Knowledge Tracing (DKT)
аргыг Bayesian Knowledge Tracing (BKT) болон логистик загваруудтай нэгдсэн магадлалын тавилд
харьцуулан шинжлэв. Загваруудыг мэдлэгийн төлөвийн дүрслэл, параметрийн тоо, чадварын
шошгоос хамаарах байдал, чадвар хоорондын шилжилт, тайлбарлах боломж гэсэн ижил шалгуураар
үнэлж, DKT-ийн хосолсон оролтын кодчилол, ганц гаралтад тооцох алдагдлын функц ба багтаамжийн
онолын үр дагаврыг гаргав. Мөн зохиогчдын нийтэлсэн ASSISTments 2009 өгөгдлийг аудит хийхэд
4\,217 сурагч, 525\,534 хариулт гарсан нь өгүүлэлд мэдээлсэн 15\,931 сурагчтай зөрөв. Нэгээс
олон чадвар бодсон 3\,398 сурагч бүгдийн хариулт чадварын дугаараар эрэмбэлэгдсэн байх ба
дараалсан хариултын 92.4\% нь ижил чадвартай байв. Ийм бүтэцтэй өгөгдөлд параметргүй
«тухайн чадварын өмнөх хариултаа давтах» дүрэм AUC 0.719 өгч, өгүүлэлд мэдээлсэн BKT-ийн
0.67-оос давсан. Эдгээр нь DKT-ийн 0.86 AUC-ийн нэг хэсэг өгөгдлийн эрэмбэ, орон нутгийн
хамаарлаас үүдэлтэй байж болзошгүйг, ялангуяа чадвар хоорондын шилжилт ба бодлогын дарааллын
дүгнэлтийг болгоомжтой тайлбарлах шаардлагатайг харуулна. Эцэст нь бага өгөгдөлтэй
платформд загвар сонгох, үнэлгээ хийх практик зөвлөмж гаргав.}
\par\vspace{5pt}
\noindent\textbf{\textit{Түлхүүр үг:}} \textit{knowledge tracing, LSTM, Bayesian knowledge tracing,
ASSISTments, AUC, өгөгдлийн аудит}
\end{minipage}
\vspace{14pt}
}]

%===============================================================================
\section{Удиртгал}

Онлайн сургалтын системүүд сурагчийн бодлого бүр дээрх оролдлогыг хадгалдаг болсноор тэдний
мэдлэгийн төлөвийг тусгай шалгалтгүйгээр, тасралтгүй үнэлэх боломж бүрдсэн. Ийм үнэлгээ нь
сурагчид тохирох хүндрэлтэй бодлого сонгох, сэдвийг эзэмшсэн эсэхийг шийдэх, сургалтын
дарааллыг хувь хүнд тааруулахад шаардлагатай. Энэ даалгаврыг \emph{knowledge tracing} (мэдлэг
мөшгих) гэдэг.

Уг салбарын анхны өргөн хэрэглэгдсэн загвар нь чадвар бүрийг хоёртын далд хувьсагчаар
илэрхийлдэг Bayesian Knowledge Tracing (BKT) юм [2]. Үүнтэй зэрэгцэн сурагчийн чадвар ба
бодлогын хүндрэлийг логистик функцээр холбодог Item Response Theory [14], түүний онлайн хувилбар
Elo систем [11], өмнөх зөв, буруу хариултын тоог шинж болгодог Performance Factors Analysis
[10] хөгжсөн. 2015 онд Piech нар рекуррент мэдрэлийн сүлжээнд суурилсан Deep Knowledge Tracing
(DKT)-ийг санал болгож, жишиг өгөгдөл дээр AUC-ийн 25\%-ийн ахиц мэдээлсэн [12]. Энэ үр дүн
дараагийн гүн сургалтын загваруудын эхлэл болсон [1], [9], [16]. Харин хожмын судалгаанууд
сайжруулсан BKT ба логистик загварууд DKT-тэй ойролцоо гүйцэтгэл үзүүлж болохыг [5], [8],
жишиг өгөгдөлд өгөгдлийн асуудал байсныг [15] мэдээлсэн.

Иймд DKT-ийн аль загварын шийдвэр үнэхээр ач холбогдолтой, жишиг өгөгдлийн аль шинж нь
мэдээлсэн тоонд нөлөөлсөн бэ гэсэн асуулт нээлттэй хэвээр байна. Энэ асуулт бага өгөгдөлтэй
шинэ платформд загвар сонгоход ч чухал. Энэ ажлын хувь нэмэр нь: (i) мэдлэг мөшгих загваруудыг
нэгдсэн «төлөв–шинэчлэл–унших» тавилд оруулж, ижил шалгуураар харьцуулах; (ii) DKT-ийн оролтын
кодчилол, алдагдлын функц, багтаамжийн онолын үр дагаврыг гаргах; (iii) зохиогчдын нийтэлсэн
ASSISTments өгөгдлийг аудит хийж, үр дүнг тайлбарлахад нөлөөлөх бүтцийн шинжийг тоогоор
харуулах; (iv) эдгээрээс бага өгөгдөлтэй системд хэрэгжих зөвлөмж гаргах явдал юм. Шинэ
загвар сургах туршилтыг энэ ажилд хийгээгүй бөгөөд дараагийн шатанд үлдээв.

%===============================================================================
\section{Бодлогын тавил ба үнэлгээ}

\subsection{Тэмдэглэгээ ба нэгдсэн тавил}
Сурагчийн $t$ дахь оролдлогыг $x_t=(q_t,a_t)$ гэж тэмдэглэнэ. Энд $q_t\in\{1,\dots,M\}$ нь
бодлогын (чадварын) таг, $a_t\in\{0,1\}$ нь зөв эсэх. Загвар нь түүх
$H_t=(x_1,\dots,x_t)$ болон дараагийн бодлого $q_{t+1}$-ээс
\begin{equation}
\hat p_{t+1}=P_\theta\bigl(a_{t+1}=1\mid q_{t+1},H_t\bigr)
\end{equation}
магадлалыг гаргана. Мэдлэг мөшгих бүх загварыг гурван бүрэлдэхүүнээр илэрхийлж болно: мэдлэгийн
төлөв $s_t$, шинэчлэл $s_t=g_\theta(s_{t-1},x_t)$, унших функц
$\hat p_{t+1}=r_\theta(s_t,q_{t+1})$. Загварууд хоорондын ялгаа нь $s_t$-ийн хэмжээс ба
$g,r$-ийн уян хатан байдалд оршино (1-р хүснэгт).

\subsection{AUC хэмжүүр ба түүний онцлог}
Мэдлэг мөшгих ажлуудын гол хэмжүүр нь ROC муруйн доорх талбай (AUC) юм. Энэ нь санамсаргүй
сонгосон зөв хариултын таамаглал санамсаргүй буруу хариултынхаас өндөр байх магадлал [6]:
\begin{equation}
\mathrm{AUC}=P\bigl(\hat p^{+}>\hat p^{-}\bigr)+\tfrac12 P\bigl(\hat p^{+}=\hat p^{-}\bigr).
\end{equation}
Энэ тодорхойлолтоос хоёр үр дагавар гарна. Нэгдүгээрт, AUC нь зөвхөн эрэмбийг хэмждэг тул
магадлалын тохируулгыг (calibration) үнэлдэггүй. Харин адаптив систем «эзэмшсэн» гэж шийдэхдээ
ихэвчлэн тогтмол босго ашигладаг тул энэ нь тусдаа шалгах ёстой шинж юм. Хоёрдугаарт, AUC-г
бүх (сурагч,~$t$) хос дээр нэгтгэж тооцдог тул урт дараалалтай цөөн сурагч үр дүнг давамгайлна.
Бидний аудитаар ASSISTments-ийн хамгийн урт дараалалтай 10\% сурагч нийт хариултын 64.8\%-ийг
бүрдүүлж байна (3-р хүснэгт).

%===============================================================================
\section{Уламжлалт загварууд}

\subsection{Bayesian Knowledge Tracing}
BKT нь чадвар бүрт $L\in\{\text{мэднэ},\text{мэдэхгүй}\}$ далд хувьсагчтай Hidden Markov
Model юм [2]. Параметр нь анхны мэдлэг $P(L_0)$, сурах $T$, таамаглах $G$, алдах $S$.
Хариулт ажиглагдсаны дараа Байесын дүрмээр
\begin{equation}
P(L_t\mid a_t{=}1)=\frac{P(L_t)(1-S)}{P(L_t)(1-S)+\bigl(1-P(L_t)\bigr)G}
\end{equation}
(буруу үед $1-S\to S$, $G\to 1-G$) шинэчилж, дараа нь сурах шилжилт хийнэ:
\begin{equation}
P(L_{t+1})=P(L_t\mid a_t)+\bigl(1-P(L_t\mid a_t)\bigr)T.
\end{equation}
Таамаглал нь $\hat p=P(L)(1-S)+(1-P(L))G$. \ref{fig:bkt}-р зурагт $P(L_0)=0.2$, $T=0.15$, $G=0.2$,
$S=0.1$ үед найман хариултын жишээг бид тооцоолж үзүүлэв. Зөв хариулт бүр $P(L)$-г огцом
өсгөж, 4 дэх алхмын буруу хариулт түүнийг 0.56-аас 0.14 руу бууруулна. Энэ нь BKT-ийн төлөв нэг
скаляр тул сүүлийн ажиглалтад хүчтэй хариу үйлдэл үзүүлдгийг харуулна. Мөн чадварууд
хоорондоо хамааралгүй тул нэг чадварын ахиц нөгөөгийн таамаглалд нөлөөлөхгүй.

\begin{figure}[t]
\centering
\begin{tikzpicture}
\begin{axis}[width=\columnwidth,height=4.2cm,xmin=0.5,xmax=8.5,ymin=0,ymax=1.0,clip=false,
  xtick={1,...,8},ytick={0,0.25,0.5,0.75,1},xlabel={Алхам $t$},
  tick label style={font=\scriptsize},label style={font=\small},
  legend style={font=\scriptsize,at={(0.98,0.04)},anchor=south east,draw=none,fill=white},
  grid=major,grid style={gray!20}]
\addplot[thick,black,mark=*,mark size=1.6pt] coordinates
  {(1,0.200)(2,0.176)(3,0.172)(4,0.561)(5,0.267)(6,0.678)(7,0.919)(8,0.984)};
\addlegendentry{$P(L_t)$}
\addplot[thick,gray,dashed,mark=square*,mark size=1.4pt] coordinates
  {(1,0.340)(2,0.323)(3,0.320)(4,0.593)(5,0.387)(6,0.675)(7,0.843)(8,0.889)};
\addlegendentry{$\hat p_t$}
\node[font=\scriptsize] at (axis cs:1,1.09) {0};
\node[font=\scriptsize] at (axis cs:2,1.09) {0};
\node[font=\scriptsize] at (axis cs:3,1.09) {1};
\node[font=\scriptsize] at (axis cs:4,1.09) {0};
\node[font=\scriptsize] at (axis cs:5,1.09) {1};
\node[font=\scriptsize] at (axis cs:6,1.09) {1};
\node[font=\scriptsize] at (axis cs:7,1.09) {1};
\node[font=\scriptsize] at (axis cs:8,1.09) {1};
\end{axis}
\end{tikzpicture}
\caption{BKT-ийн жишээ ($P(L_0)=0.2$, $T=0.15$, $G=0.2$, $S=0.1$). Дээд мөрөнд
ажиглагдсан хариулт $a_t$; $\hat p_t$ нь тухайн алхмын өмнөх таамаглал (бидний тооцоо).}
\label{fig:bkt}
\end{figure}

\subsection{Логистик загварууд}
Rasch (1PL IRT) загвар нь $\operatorname{logit}\hat p=\theta_s-\beta_q$ хэлбэртэй бөгөөд
сурагчийн чадвар $\theta_s$, бодлогын хүндрэл $\beta_q$-гээр илэрхийлэгдэнэ [14]. Elo систем
нь үүнийг онлайнаар $\theta_s\leftarrow\theta_s+K(a-\hat p)$ гэж шинэчилдэг тул бага өгөгдөлд
тогтвортой ажилладаг [11]. PFA нь $\operatorname{logit}\hat p=\beta_q+\gamma_q s_q+\rho_q f_q$
гэж чадвар бүрийн өмнөх зөв $s_q$, буруу $f_q$ хариултын тоог ашиглана [10].

\begin{table*}[t]
\tabtitle{Мэдлэг мөшгих загваруудын нэгдсэн тавил дахь харьцуулалт}
\centering\small\setlength{\tabcolsep}{4pt}
\begin{tabular}{|L{1.8cm}|L{3.0cm}|L{3.8cm}|L{2.4cm}|L{2.8cm}|L{1.9cm}|}
\hline
\textbf{Загвар} & \textbf{Төлөв $s_t$} & \textbf{Параметрийн тоо} &
\textbf{Чадварын шошго} & \textbf{Чадвар хоорондын шилжилт} & \textbf{Тайлбарлах}\\\hline
BKT [2] & чадвар бүрт $P(L)\in[0,1]$ & $4M$ & шаардана & байхгүй & өндөр\\\hline
IRT/Elo [11], [14] & сурагчийн нэг $\theta_s$ & $M$ + сурагчийн тоо & бодлогоор хангалттай &
нэг $\theta_s$-ээр далд & өндөр\\\hline
PFA [10] & тоолуур $s_q,f_q$ & $3M$ & шаардана & байхгүй & өндөр\\\hline
DKT [12] & $\h_t\in\mathbb R^{H}$ & $4H(2M{+}H{+}1)+M(H{+}1)$ & заавал биш &
өгөгдлөөс сурна & бага\\\hline
\end{tabular}
\end{table*}

%===============================================================================
\section{DKT-ийн онолын шинжилгээ}

\subsection{Бүтэц}
DKT нь $(q_t,a_t)$ хосыг $2M$ урттай one-hot вектор $\x_t$ болгож LSTM [7]-д оруулна:
\begin{equation}
(\h_t,\mathbf c_t)=\mathrm{LSTM}(\x_t,\h_{t-1},\mathbf c_{t-1}),\quad
\y_t=\sigma\bigl(W\h_t+\mathbf b\bigr).
\end{equation}
$\y_t\in(0,1)^M$ нь бүх бодлогын таамаглал бөгөөд дараагийн бодлогын магадлалыг
$\hat p_{t+1}=\y_t^{\top}\delta(q_{t+1})$ гэж уншина (\ref{fig:dkt}-р зураг). Өгүүлэлд далд давхаргын хэмжээ
200, минибатч 100, гаралтад л хэрэглэх dropout, градиент таслалт ашигласан [12].

\begin{figure}[t]
\centering
\begin{tikzpicture}[node distance=4mm and 5.5mm, every node/.style={font=\scriptsize},
  box/.style={draw,rounded corners,minimum width=10mm,minimum height=5mm},
  arr/.style={-{Stealth[length=1.6mm]}}]
  \node[box,fill=gray!10] (h1) {LSTM};
  \node[box,fill=gray!10,right=of h1] (h2) {LSTM};
  \node[box,fill=gray!10,right=of h2] (h3) {LSTM};
  \node[right=3mm of h3] (dots) {$\cdots$};
  \node[left=3mm of h1] (h0) {$\h_0$};
  \foreach \i in {1,2,3}{
    \node[below=of h\i,draw,fill=gray!12] (x\i) {$(q_\i,a_\i)$};
    \node[above=of h\i,draw,fill=gray!4] (y\i) {$\y_\i$};
    \draw[arr] (x\i)--(h\i); \draw[arr] (h\i)--(y\i);}
  \draw[arr] (h0)--(h1); \draw[arr] (h1)--node[above]{$\h_1$}(h2);
  \draw[arr] (h2)--node[above]{$\h_2$}(h3); \draw[arr] (h3)--(dots);
\end{tikzpicture}
\caption{DKT-ийн хугацаагаар задалсан бүтэц. Алхам бүрт $\y_t$-ийн зөвхөн $q_{t+1}$
элемент алдагдалд ордог. Оролт нь $(q_t,a_t)$ хосын one-hot $e_{q_t+Ma_t}$.}
\label{fig:dkt}
\end{figure}

\subsection{Хосолсон one-hot кодчилол}
Зохиогчид $q_t$, $a_t$-г тусад нь кодлоход гүйцэтгэл муудсан гэж тэмдэглэсэн боловч шалтгааныг
тайлбарлаагүй [12]. Бид үүнийг эхний давхаргын шугаман хэлбэрээр тайлбарлаж байна. Тусдаа
кодчилол $\x_t=[e_{q_t};a_t]$ үед LSTM-ийн оролтын урьдчилсан идэвхжүүлэлт
\begin{equation}
W_x\x_t=W^{(q)}_{:,q_t}+a_t\,\mathbf w^{(a)}
\end{equation}
болно. Зөв хариултын нөлөө $\mathbf w^{(a)}$ нь бүх бодлогод \emph{ижил вектор} тул «$q_1$-г зөв
бодсон», «$q_2$-г зөв бодсон» гэсэн ялгааг эхний давхарга нэмэх хэлбэрээр л илэрхийлнэ. Бодлого
ба үр дүнгийн харилцан үйлчлэлийг рекуррент шугаман бус хувиргалт дараа нь сурах шаардлагатай.
Хосолсон кодчилол $\x_t=e_{q_t+Ma_t}$ үед $W_x\x_t=W_{:,\,q_t+Ma_t}$ буюу $(q,a)$ хос бүр
өөрийн баганатай болж, харилцан үйлчлэл эхний давхаргад шууд бий болно. Үнэ нь оролтын жингийн
тоо $(M{+}1)H$-аас $2MH$ болж хоёр дахин өсөх явдал юм.

Бодлогын тоо их үед зохиогчид $(q,a)$ хос бүрт $N\ll 2M$ хэмжээст санамсаргүй Гаусс вектор оноохыг
санал болгож, үүнийг compressed sensing-ээр үндэслэсэн. One-hot векторууд хоорондоо ижил
$\sqrt2$ зайтай тул энэ нь Johnson–Lindenstrauss леммийн [3] нөхцөлд ч тайлбарлагдана:
$N=O(\varepsilon^{-2}\log 2M)$ хэмжээс нь бүх хос зайг $1\pm\varepsilon$ нарийвчлалтай хадгална.

\subsection{Ганц гаралтын хяналт}
Алдагдлын функц нь $\mathcal L=\sum_t\ell(\y_t^{\top}\delta(q_{t+1}),a_{t+1})$, энд $\ell$ нь
binary cross-entropy [12]. Гаралтын logit $z_{t,j}$-ээр уламжлал авбал
\begin{equation}
\frac{\partial\mathcal L}{\partial z_{t,j}}=\bigl(\sigma(z_{t,j})-a_{t+1}\bigr)\,
\mathbb 1[j=q_{t+1}]
\end{equation}
болно. Өөрөөр хэлбэл алхам бүрт $W$-ийн зөвхөн нэг мөр градиент авч, бусад $M-1$ таамаглалд
хяналт байхгүй. Үүнээс хоёр үр дагавар гарна: (a) одоогийн бодлогыг зөв бодсон ч тухайн бодлогын
таамаглал өсөх баталгаагүй; (b) ажиглагдаагүй бодлогуудын таамаглал алхам алхмаар хэлбэлзэж
болно. Yeung, Yeung эдгээрийг одоогийн хариултыг сэргээх ба $\y_t-\y_{t-1}$-ийг жигдрүүлэх
нэмэлт алдагдлаар засахыг санал болгосон [16].

\subsection{Багтаамж ба өгөгдлийн хэрэгцээ}
1-р хүснэгтийн томьёогоор ASSISTments ($M=124$, $H=200$) дээр DKT нь
$359\,200+24\,924=384\,124$ параметртэй (PyTorch хоёр bias ашиглавал 384\,924). Энэ нь BKT-ийн
496 параметрээс 774 дахин их бөгөөд нэг параметрт ердөө 1.37 хариулт ногдоно. Иймд DKT-ийн
давуу тал өгөгдлийн хэмжээнээс хүчтэй хамаарах нь гайхалтай биш. Gervet нар олон өгөгдөл дээр
харьцуулж, бага өгөгдөлд логистик загварууд өрсөлдөхүйц, их өгөгдөлд гүн загварууд давуу
болохыг мэдээлсэн [5].

%===============================================================================
\section{Жишиг өгөгдлийн аудит}

\subsection{Арга зүй}
Өгүүллийн мэдээлсэн үр дүнг (2-р хүснэгт) тайлбарлахын тулд зохиогчдын GitHub репод нийтэлсэн
ASSISTments 2009--2010 «skill builder» [4] өгөгдлийн \texttt{builder\_train.csv},
\texttt{builder\_test.csv} файлыг [13] өөрчлөлтгүйгээр уншиж, энгийн тоолох скриптээр
статистик гаргав. Файлд сурагч бүрийн хариултын тоо, чадварын таг, зөв эсэхийн дараалал
хадгалагдсан. Ижил чадварын тасралтгүй хэсгийг «блок» гэж нэрлэв.

\begin{table}[t]
\tabtitle{Өгүүлэлд мэдээлсэн AUC [12]}
\centering\small
\setlength{\tabcolsep}{4pt}
\begin{tabular}{|l|c|c|c|c|}
\hline
\textbf{Өгөгдөл} & \textbf{Marginal} & \textbf{BKT} & \textbf{BKT*} & \textbf{DKT}\\\hline
Simulated-5 & -- & 0.54 & -- & 0.75\\\hline
Khan Math & 0.63 & 0.68 & -- & 0.85\\\hline
ASSISTments & 0.62 & 0.67 & 0.69 & 0.86\\\hline
\end{tabular}
\end{table}

\subsection{Үр дүн}
3-р хүснэгтэд аудитын үр дүнг нэгтгэв. Хариултын тоо (525\,534) ба таг (124) өгүүлэлтэй
таарч байгаа ч сурагчийн тоо 4\,217 бөгөөд өгүүлэлд мэдээлсэн 15\,931-ээс 3.8 дахин бага байна.
Блокийн тоо (44\,087) ч энэ тоотой таарахгүй тул зөрүүний шалтгааныг файлаас тогтоох боломжгүй.

\begin{table}[t]
\tabtitle{Нийтэлсэн ASSISTments файлын аудит (бидний тооцоо)}
\centering\small
\setlength{\tabcolsep}{3pt}
\begin{tabular}{|L{5.2cm}|r|}
\hline
\textbf{Үзүүлэлт} & \textbf{Утга}\\\hline
Сурагч (өгүүлэлд 15\,931) & 4\,217\\\hline
Хариулт (өгүүлэлд 526K) / таг & 525\,534 / 124\\\hline
Дарааллын урт: медиан / дундаж / их & 27 / 124.6 / 8\,214\\\hline
10-аас цөөн хариулттай сурагч & 23.1\%\\\hline
Хамгийн урт 10\% сурагчийн хариултын хувь & 64.8\%\\\hline
Чадварын дугаараар эрэмбэлэгдсэн сурагч & 3\,398/3\,398\\\hline
Өмнөхтэйгээ ижил чадвартай хариулт & 92.4\%\\\hline
Блок: тоо / дундаж урт & 44\,087 / 11.9\\\hline
$P(a_t{=}1\mid a_{t-1}{=}1)$, блок дотор & 0.828\\\hline
$P(a_t{=}1\mid a_{t-1}{=}0)$, блок дотор & 0.407\\\hline
«Өмнөх хариултаа давтах» дүрмийн AUC & 0.719\\\hline
\end{tabular}
\end{table}

Хамгийн чухал олдвор нь дарааллын эрэмбэ юм. Нэгээс олон чадвар бодсон 3\,398 сурагч
бүгдийн хариулт чадварын дугаарын өсөх эрэмбээр байрласан бөгөөд нэг сурагчийн дараалалд аль
ч чадвар хоёр дахь удаа гарч ирэхгүй. Иймд файл дахь дараалал чадвар хооронд
\emph{хугацааны} эрэмбэ биш, харин чадварын дугаарын эрэмбийг илэрхийлнэ. Эх өгөгдөл чадвараар
эрэмбэлэгдсэн байсан нь үүний шалтгаан байж болзошгүй ч бид эх өгөгдлийг шалгаагүй.

Блок дотор хариултууд хүчтэй хамааралтай: өмнөх хариулт зөв бол дараагийнх зөв байх магадлал
0.828, буруу бол 0.407. Блок дахь байрлалаар зөв хувь 0.581-ээс 0.70 хүртэл өсдөг (\ref{fig:block}-р зураг).
Үүнийг ашиглан бид сургалтгүй, параметргүй дүрэм туршив: тухайн чадварын өмнөх хариулт
байвал түүнийг, байхгүй бол 0.5-ыг таамаглал болгоно. Энэ дүрэм 521\,317 таамаглал дээр AUC
0.719 өгч, өгүүлэлд мэдээлсэн BKT (0.67), BKT* (0.69)-ээс давсан. Дүрэм нь сургагдах параметргүй
тул өгөгдөл алдагдах (leakage) асуудалгүй.

\begin{figure}[t]
\centering
\begin{tikzpicture}
\begin{axis}[width=\columnwidth,height=3.9cm,ybar,bar width=8pt,ymin=0.5,ymax=0.74,
  xtick={1,...,10},xticklabels={1,2,3,4,5,6,7,8,9,${\ge}10$},
  xlabel={Блок дахь байрлал},ylabel={Зөв хувь},
  tick label style={font=\scriptsize},label style={font=\small},
  grid=major,grid style={gray!20}]
\addplot[fill=gray!35,draw=black] coordinates
 {(1,0.581)(2,0.653)(3,0.665)(4,0.687)(5,0.682)(6,0.675)(7,0.673)(8,0.681)(9,0.675)(10,0.699)};
\end{axis}
\end{tikzpicture}
\caption{Ижил чадварын блок дахь байрлалаар зөв хариултын хувь (нийтэлсэн ASSISTments файл,
бидний тооцоо).}
\label{fig:block}
\end{figure}

\subsection{Үр дүнгийн тайлбарт үзүүлэх нөлөө}
Нэгдүгээрт, дараагийн хариултын 92.4\% нь ижил чадварынх, тэдгээр нь хоорондоо хүчтэй
хамааралтай тул сүүлийн хариултыг хадгалах чадвартай аливаа загвар өндөр AUC авна. Иймд
0.86-ийн нэг хэсэг нь чадвар хоорондын нарийн бүтэц биш, орон нутгийн хамаарлаас үүдэлтэй байж
болзошгүй. Энэ нь сайжруулсан BKT DKT-тэй ойртсон [8], засварласан өгөгдөл дээр DKT-ийн AUC буурсан
[15] гэсэн мэдээлэлтэй нийцнэ.

Хоёрдугаарт, чадвар хоорондын шилжилт файлд чадварын дугаарын эрэмбээр л явагддаг. Иймд энэ
файлаар сургасан загварын чадвар хоорондын нөлөө ($J_{ij}$) нь бодит хугацааны сургалтын
дарааллыг бус, дугаарын эрэмбийг тусгах эрсдэлтэй. Мөн ийм өгөгдлөөр сургасан загвар зөвхөн
блоклосон дараалал харсан тул сэдвийг холих (mixing) дараалал түүний хувьд сургалтын тархалтаас
гадуурх оролт болно. ASSISTments дээрх «blocking нь mixing-ээс давуу» гэсэн загвар дээрх
симуляцийн дүгнэлт [12] үүнээс хэсэгчлэн хамаарч болзошгүй.

Гуравдугаарт, DKT-ийн нэхэмжлэлүүд өөр өөр түвшний нотолгоонд тулгуурладаг (4-р хүснэгт).
Тэдгээрийг ижил хүчтэй баримт мэт авч үзэх нь буруу.

\begin{table}[t]
\tabtitle{DKT-ийн нэхэмжлэл ба нотолгооны төрөл}
\centering\small
\begin{tabular}{|L{2.2cm}|L{1.8cm}|L{3.0cm}|}
\hline
\textbf{Нэхэмжлэл} & \textbf{Нотолгоо} & \textbf{Анхаарах зүйл}\\\hline
AUC-ийн 25\% ахиц & Жишиг өгөгдөл & Өгөгдлийн хувилбар, эрэмбэ, хуваалтаас хамаарна\\\hline
Шошгогүй далд бүтэц сурна & Симуляц & IRT-ээр үүсгэсэн өгөгдөл\\\hline
Урьдчилсан нөхцөлийн граф & Чанарын дүрслэл & Шинжээчийн графтай тоон харьцуулалтгүй\\\hline
Blocking > mixing & Загвар дээрх симуляц & Бодит сурагчтай туршилт биш\\\hline
\end{tabular}
\end{table}

%===============================================================================
\section{Хэлэлцүүлэг}

DKT-ийн бодит хувь нэмэр нь мэдлэгийн төлөвийг шинжээчийн шошгогүйгээр, олон хэмжээст векторт
сурах боломжийг нээж, дараагийн attention загваруудын үндсийг тавьсан явдал юм [1], [9]. Харин
бидний шинжилгээ DKT-ийн давуу талыг гурван нөхцөлтэйгөөр ойлгох ёстойг харуулж байна: (a) хангалттай
их өгөгдөл, (b) хугацааны зөв эрэмбэтэй дараалал, (c) орон нутгийн хамаарлыг аль хэдийн
ашигладаг хүчтэй суурь загвартай харьцуулалт.

Эдгээрээс бага өгөгдөлтэй шинэ платформд, жишээлбэл ерөнхий боловсролын математикийн бодлогын
сайтад дараах зөвлөмж гарна. (1) Оролдлогыг цаг хугацааны тэмдэгтэй, жинхэнэ дарааллаар нь
хадгалах. (2) Эхний шатанд Elo, IRT, PFA зэрэг цөөн параметртэй загвараас эхлэх. Тэдгээр нь
тайлбарлах боломжтой бөгөөд бага өгөгдөлд тогтвортой. (3) Үнэлгээг сурагчаар хуваасан, хугацааны
эрэмбийг хадгалсан хуваалтаар хийж, AUC-ээс гадна тохируулгын хэмжүүр тайлагнах. (4) «Өмнөх хариултаа
давтах» шиг энгийн дүрмийг заавал суурь загварт оруулах. (5) Хариултын тоо параметрийн тооноос
олон дахин их болсны дараа DKT руу шилжих.

Энэ ажлын хязгаарлалт: аудит нь зөвхөн нийтэлсэн файлын тоолох шинжилгээ бөгөөд эх ASSISTments
өгөгдөл, Khan өгөгдөлд хандаагүй. Мөн DKT-ийг дахин сургаж, эрэмбийг сэргээсэн өгөгдөл дээрх
нөлөөг хэмжээгүй. Энэ нь дараагийн туршилтын гол зорилт болно.

%===============================================================================
\section*{Дүгнэлт}

DKT нь мэдлэг мөшгихийг LSTM-ээр шийдэж, хосолсон one-hot оролт ба зөвхөн дараагийн бодлогод
тооцох алдагдлаар шинжээчийн шошгогүйгээр өндөр AUC гаргасан. Бидний онолын шинжилгээ хосолсон
кодчилол эхний давхаргад бодлого ба үр дүнгийн харилцан үйлчлэлийг шууд үүсгэдгийг, ганц
гаралтын хяналт таамаглалын тогтворгүй байдлын эх үүсвэр болохыг, 384 мянган параметр нь BKT-ээс
774 дахин их өгөгдлийн хэрэгцээ үүсгэдгийг харуулав. Нийтэлсэн ASSISTments файлын аудит нь сурагчийн
тооны зөрүү, чадварын дугаараар эрэмбэлэгдсэн дараалал, блок доторх хүчтэй хамаарлыг илрүүлж,
параметргүй дүрэм хүртэл AUC 0.719 авдгийг харуулав. Иймд DKT-ийн мэдээлсэн ахицыг өгөгдлийн
эрэмбэ ба суурь загварын хүчнээс салгаж үнэлэх шаардлагатай. Бага өгөгдөлтэй системд цөөн
параметртэй загвараас эхэлж, хугацааны эрэмбийг хадгалсан өгөгдөл хуримтлуулах нь зөв алхам юм.

%===============================================================================
\renewcommand{\refname}{\normalsize\centering\bfseries АШИГЛАСАН МАТЕРИАЛ, НОМ ЗҮЙ}
\begin{thebibliography}{99}\footnotesize
\setlength{\itemsep}{1pt}
\bibitem{abd} G. Abdelrahman, Q. Wang, and B. Nunes, “Knowledge tracing: A survey,”
\emph{ACM Comput. Surv.}, vol. 55, no. 11, 2023.
\bibitem{corbett} A. T. Corbett and J. R. Anderson, “Knowledge tracing: Modeling the acquisition of
procedural knowledge,” \emph{User Model. User-Adapt. Interact.}, vol. 4, no. 4, pp. 253--278, 1994.
\bibitem{jl} S. Dasgupta and A. Gupta, “An elementary proof of a theorem of Johnson and
Lindenstrauss,” \emph{Random Struct. Algorithms}, vol. 22, no. 1, pp. 60--65, 2003.
\bibitem{feng} M. Feng, N. Heffernan, and K. Koedinger, “Addressing the assessment challenge with
an online system that tutors as it assesses,” \emph{User Model. User-Adapt. Interact.}, vol. 19,
no. 3, pp. 243--266, 2009.
\bibitem{gervet} T. Gervet, K. Koedinger, J. Schneider, and T. Mitchell, “When is deep learning the
best approach to knowledge tracing?” \emph{J. Educ. Data Min.}, vol. 12, no. 3, pp. 31--54, 2020.
\bibitem{hanley} J. A. Hanley and B. J. McNeil, “The meaning and use of the area under a receiver
operating characteristic (ROC) curve,” \emph{Radiology}, vol. 143, no. 1, pp. 29--36, 1982.
\bibitem{lstm} S. Hochreiter and J. Schmidhuber, “Long short-term memory,” \emph{Neural Comput.},
vol. 9, no. 8, pp. 1735--1780, 1997.
\bibitem{khajah} M. Khajah, R. V. Lindsey, and M. C. Mozer, “How deep is knowledge tracing?” in
\emph{Proc. 9th Int. Conf. Educ. Data Min. (EDM)}, 2016, pp. 94--101.
\bibitem{sakt} S. Pandey and G. Karypis, “A self-attentive model for knowledge tracing,” in
\emph{Proc. 12th Int. Conf. Educ. Data Min. (EDM)}, 2019, pp. 384--389.
\bibitem{pfa} P. I. Pavlik Jr., H. Cen, and K. R. Koedinger, “Performance factors analysis---A new
alternative to knowledge tracing,” in \emph{Proc. 14th Int. Conf. Artif. Intell. Educ. (AIED)},
2009, pp. 531--538.
\bibitem{elo} R. Pelánek, “Applications of the Elo rating system in adaptive educational
systems,” \emph{Comput. Educ.}, vol. 98, pp. 169--179, 2016.
\bibitem{piech} C. Piech, J. Bassen, J. Huang, S. Ganguli, M. Sahami, L. Guibas, and
J. Sohl-Dickstein, “Deep knowledge tracing,” in \emph{Adv. Neural Inf. Process. Syst. 28
(NeurIPS)}, 2015, pp. 505--513.
\bibitem{repo} C. Piech \emph{et al.}, “DeepKnowledgeTracing,” GitHub repository,
github.com/chrispiech/DeepKnowledgeTracing, accessed Nov. 19, 2026.
\bibitem{rasch} G. Rasch, \emph{Probabilistic Models for Some Intelligence and Attainment Tests}.
Copenhagen, Denmark: Danish Institute for Educational Research, 1960.
\bibitem{xiong} X. Xiong, S. Zhao, E. G. Van Inwegen, and J. E. Beck, “Going deeper with deep
knowledge tracing,” in \emph{Proc. 9th Int. Conf. Educ. Data Min. (EDM)}, 2016, pp. 545--550.
\bibitem{dktp} C.-K. Yeung and D.-Y. Yeung, “Addressing two problems in deep knowledge tracing
via prediction-consistent regularization,” in \emph{Proc. 5th ACM Conf. Learning @ Scale (L@S)},
2018, Art. no. 5.
\end{thebibliography}

\end{document}
